% FILE: 02_lineage_and_context.tex
% Project: otel-weaver-exp-002-diff-to-residual
% Thesis Role: weaver-diff-to-pi-residual-bridge-experiment

\section{Lineage and Context}
\label{sec:otel-weaver-exp-002-diff-to-residual-lineage}

\subsection{2016 Language-Model Lesson}
\label{sec:otel-weaver-exp-002-diff-to-residual-lineage-2016}

The 2016 language-model experiment established a foundational finding that is
directly ancestral to EXP-002: local text imitation does not conserve consequence.
A model can reproduce the surface structure of an utterance---its vocabulary,
syntax, and statistical patterns---without capturing the causal chain that gives
the utterance its meaning in a process context. The model ``knows'' the words; it
does not know what must have happened before those words could be lawfully spoken.

EXP-002 is the 2016 lesson applied to telemetry engineering. OTel Weaver is, in
this analogy, a surface-structure validator: it confirms that a span contains the
expected field names and value types. It is a sophisticated schema-imitation check.
What it cannot do---and was never designed to do---is assert that the activity
those fields describe occurred in a lawful sequence relative to other activities.
The doctrine file \texttt{otel-weaver-is-feedstock.md} states this precisely:
``Structural alignment is merely the verification that our instruments are recording
within expected parameters. It is an engineering metric.''

The schema rename from \texttt{process.pi.activity.name} to
\texttt{process.pi.transition.name} is a surface-structure change. It changes the
labels on the instruments without altering a single execution path in the underlying
process. The 2016 lesson predicts that any system which treats surface structure as
consequence will misread this rename as a consequence-level event---a process
violation. EXP-002 manufactures the architectural mechanism that prevents this
misreading by ensuring surface-structure changes are resolved before they reach the
consequence-evaluation layer.

\subsection{Enterprise Process Gap}
\label{sec:otel-weaver-exp-002-diff-to-residual-lineage-enterprise-gap}

The enterprise process gap that motivates the broader research program is the absence
of a full-lifecycle process intelligence stack in commercial observability tooling.
Existing telemetry platforms---including those built on OTel Weaver---are designed for
infrastructure monitoring: latency histograms, error rates, service dependency maps.
They are not designed to answer the question: ``Did this business process execute
in accordance with its normative model, and can that execution be proven to an
auditor?''

EXP-002 addresses a specific manifestation of this gap: when a telemetry schema is
updated during normal platform maintenance, commercial observability tools will either
silently accept the new schema (losing conformance signal) or flag every new-schema
record as malformed (generating false compliance alerts). Neither response is lawful
under a full-lifecycle process intelligence doctrine.

The gap is structural rather than incidental. It arises because the telemetry layer
and the conformance layer are not architecturally separated in commercial stacks. The
\texttt{BridgeRx} pattern is an explicit proposal for how to close this gap at the
ingestion boundary: place a typed, testable, auditable translation layer between the
feedstock plane and the court plane, so that schema evolution in one plane never
contaminates the other. The YAML boundary rules in
\texttt{mappings/registry-diff-to-pi-drift-map.yaml} formalize this as
\texttt{drift\_isolation\_rule\_1} and \texttt{drift\_isolation\_rule\_2}, routing
schema-diff events to a Weaver Diff Log rather than to the process drift court.

\subsection{Relationship to the Chatman Equation}
\label{sec:otel-weaver-exp-002-diff-to-residual-lineage-chatman-equation}

The Chatman Equation $A = \mu(O^*)$ asserts that organizational authority is a
function of the observable state of the organization's process artifacts. Authority
can only be computed over lawful, receipt-backed process evidence; it cannot be
computed over raw telemetry. \textsc{[AUTHOR THESIS / INTERPRETATION]}

EXP-002's contribution to this equation is categorical: it ensures that what enters
the $O^*$ observable surface is process evidence, not schema noise. If the
\texttt{BridgeRx} translation layer is absent, a schema rename event contaminates
the observable surface with false conformance failures, causing $\mu$ to compute
a degraded authority value from a polluted evidence set. The bridge enforces the
purity of the observable surface by resolving schema-layer noise before it can
reach the authority computation layer.

Put differently: the Chatman Equation assumes that its input $O^*$ contains only
lawful process evidence. EXP-002 is one of the engineering mechanisms that maintains
that assumption under real-world schema evolution pressure. Every schema version
transition that \texttt{BridgeRx} handles correctly is a boundary enforcement that
keeps $O^*$ clean.

\subsection{Relationship to Receipts and Replay}
\label{sec:otel-weaver-exp-002-diff-to-residual-lineage-receipts}

The receipt-and-replay pattern, central to the research program's verifiability
doctrine, requires that every process state transition be cryptographically
sealed and replayable. EXP-002 interacts with this pattern at the field-preservation
level: the \texttt{BridgeRx} \texttt{translate()} method explicitly preserves the
\texttt{witness\_hash} field through the S2-to-S1 translation.

The \texttt{src/bridge.rs} implementation assigns
\texttt{witness\_hash: input.witness\_hash} directly, with no transformation. This
ensures that the BLAKE3 cryptographic hash emitted by the instrumentation layer
survives the schema normalization step intact. A conformance court auditing the
translated record can verify the BLAKE3 hash against the original emission without
re-deriving it through the translation path.

The inline test \texttt{test\_bridge\_translation\_equivalence} asserts this
preservation explicitly:
\begin{verbatim}
assert_eq!(s1_telemetry.witness_hash, "blake3_signature_data_here");
\end{verbatim}
This assertion treats witness hash preservation as a first-class correctness
property of the bridge, not an incidental side effect. The design decision reflects
the broader research program's principle that audit integrity must be maintained
across every transformation boundary, even transformations that exist only to
resolve schema naming conventions.

The connection to \texttt{Admission<T,W>} and \texttt{Refusal} in the upstream
type-law patterns (documented in exp-003 and exp-004) means that a record
successfully translated by \texttt{BridgeRx} carries a preserved witness hash that
can be checked at the \texttt{Admission} gate, ensuring no hash forgery occurred
between emission and court evaluation.
